
\section{Related Work}
\label{sec:related}

\subsection{Bug Report Quality Assessment and Generation}

Early work on bug report quality used heuristic rules and
machine learning to detect incomplete or ambiguous reports.
Chaparro et al. assessed the quality of Steps to Reproduce
(S2R) using grammar parsing and neural sequence
labelling~\cite{chaparro2019assessing}. Zhang et al.
proposed CTQRS~\cite{zhang2022ctqrs}, a rule-based framework
that scores reports across morphological, relational, and
analytical dimensions. These studies established the
evaluation criteria that later LLM-based approaches inherited.

More recent work has shifted from detecting quality issues to
automatically fixing them using LLMs. Bo et al.~\cite{bo2024chatbr}
(ChatBR) combined a BERT-based detector with GPT-3.5-turbo to
generate missing information, achieving precision improvements
of 25.38--29.20\% over prior methods on six OSS projects.
Akyol et al.~\cite{akyol2026improbr} (ImproBR) used a
retrieval-augmented pipeline with GPT-4o mini to improve
Mojira/Minecraft reports, increasing the proportion of
complete reports from 7.9\% to 96.4\%. These approaches
demonstrate the effectiveness of proprietary LLMs in
few-shot or RAG settings, but rely on commercial APIs
that impose cost and privacy constraints.

\subsection{Open-Source LLMs with Training Adaptation}

To address cost and privacy limitations of proprietary
models, several studies have fine-tuned open-source LLMs
for bug report generation.
Acharya and Ginde~\cite{acharya2025bugquality} instruction
fine-tuned Qwen2.5-7B, Mistral-7B, and Llama3.2-3B on the
Bugzilla dataset (N=3,903) using the Alpaca-LoRA template,
with GPT-4o 3-shot as the reference baseline. Their
fine-tuned Qwen2.5-7B achieved CTQRS=77\%, ROUGE-1=0.61,
SBERT=0.85 versus GPT-4o 3-shot at CTQRS=75\%,
ROUGE-1=0.47, SBERT=0.82.
Choi and Yang~\cite{choi2025agentreport} (AgentReport)
applied QLoRA-4bit with a multi-agent pipeline on the same
Bugzilla corpus, reaching CTQRS=80.5\%.
Yang~\cite{yang2025ctqrsrl} further trained Qwen2.5-7B with
PPO reinforcement learning using CTQRS as the reward signal,
improving CTQRS from 0.46 to 0.76.

A consistent finding across these studies is that training
adaptation --- whether supervised fine-tuning or reinforcement
learning --- is necessary to achieve competitive SBERT scores,
as untrained models produce outputs with different stylistic
characteristics than the Bugzilla reference corpus.
Dın\c{c} et al.~\cite{dinc2025quality} corroborate this:
on a 10,000-report Mozilla Firefox corpus, a fine-tuned
RoBERTa model (F1=0.909) outperformed zero-shot and
few-shot LLM approaches, highlighting the value of
domain-specific training.

\subsection{Gap and Positioning}

Across all 13 papers we reviewed, Qwen2.5-7B appears
exclusively in fine-tuned or RL-adapted
configurations~\cite{acharya2025bugquality,
choi2025agentreport, yang2025ctqrsrl}. GPT-4o and
GPT-4o mini serve as few-shot baselines in several
studies~\cite{acharya2025bugquality, akyol2026improbr},
but no study has directly evaluated whether
\emph{Qwen2.5-7B in a few-shot-only setting} can
match the quality thresholds that GPT-4o achieves
under the same prompting conditions. This gap motivates
our study: if Qwen2.5-7B few-shot can replicate
GPT-4o's few-shot performance, organizations could
deploy a capable open-source model without any training
investment.
